\documentclass[a4paper,12pt]{article}
\usepackage[utf8]{inputenc}
\usepackage[T1]{fontenc}
\usepackage[french]{babel}
\usepackage{amsmath, amssymb, amsfonts}
\usepackage{array, longtable, booktabs}
\usepackage[table]{xcolor} % option table pour \rowcolors
\usepackage{geometry}
\geometry{left=1.8cm, right=1.8cm, top=2cm, bottom=2cm}
\usepackage{enumitem}
\usepackage{hyperref}
\hypersetup{colorlinks=true, linkcolor=blue, urlcolor=blue}

% Commandes pour les coches
\newcommand{\fait}{\ensuremath{\checkmark}}
\newcommand{\revoit}{\ensuremath{\circlearrowright}}
\newcommand{\casevide}{\ensuremath{\square}}

% Types de colonnes pour un contrôle précis des largeurs
\newcolumntype{C}[1]{>{\centering\arraybackslash}p{#1}}
\newcolumntype{L}[1]{>{\raggedright\arraybackslash}p{#1}}
\newcolumntype{R}[1]{>{\raggedleft\arraybackslash}p{#1}}

\title{Programmation annuelle de Mathématiques -- Terminale technologique (tronc commun) \\
	Année 2026--2027 -- Zone C (3 h/semaine)}
\author{}
\date{Version du \today}

\begin{document}
	
	\maketitle
	
	\section*{Présentation générale}
	Ce document propose une organisation spiralaire de l’enseignement des mathématiques en terminale technologique (tronc commun), conforme au programme officiel. L’année comporte 31 semaines de cours effectifs (hors vacances et jours fériés). Chaque notion est introduite une première fois, puis réinvestie et approfondie dans un second temps, à l’occasion de thèmes d’étude différents. Cette approche favorise une acquisition progressive et durable des connaissances.
	
	\medskip
	\textbf{Place de Python :} La programmation est intégrée de façon transverse et évaluée dans \emph{tous} les sujets (formatives, contrôles, DM, PBL). Les compétences algorithmiques (boucles, tests, fonctions, simulation, calcul approché, manipulation de listes) sont développées en parallèle des mathématiques et réactivées régulièrement selon le même principe spiralé.
	
	\medskip
	\textbf{Utilisation de l’IA dans les DM :} Dans chaque devoir maison, l’élève est invité à utiliser un outil d’intelligence artificielle générative (ChatGPT, Mistral, Copilot, etc.) pour :
	\begin{itemize}
		\item obtenir des explications ou des exemples supplémentaires sur une notion,
		\item suggérer des pistes de résolution,
		\item débuguer ou améliorer un code Python,
		\item confronter différentes méthodes de résolution.
	\end{itemize}
	\emph{L’IA ne doit cependant pas rédiger la solution finale à la place de l’élève.} Le raisonnement personnel et la rédaction restent exigibles. Un \textbf{journal de bord} (questions posées, réponses obtenues, critiques et adaptations apportées) est à annexer à chaque DM. Cette pratique développe l’esprit critique et la capacité à évaluer la pertinence des réponses générées.
	
	\section{Tableau récapitulatif des contenus à enseigner}
	Le tableau ci-dessous reprend l’intégralité des attendus du programme. La colonne \textbf{« Traité »} indique si le contenu a été abordé en classe (coche \fait), et la colonne \textbf{« Reprise »} signale les notions qui feront l’objet d’un second passage (\revoit) pour consolider la compréhension.
	
	\rowcolors{2}{gray!8}{white}
	\begin{longtable}{|L{14cm}|C{1.5cm}|C{1.5cm}|}
		\hline
		\textbf{Thème / Contenu} & \textbf{Traité} & \textbf{Reprise} \\
		\hline
		\endhead
		\hline
		\multicolumn{3}{r}{Suite de la page suivante} \\
		\endfoot
		\hline
		\endlastfoot
		
		% --- Analyse – Suites numériques ---
		\multicolumn{3}{|L{14cm}|}{\textbf{Analyse – Suites numériques}} \\
		Suites arithmétiques : moyenne, expression du terme général, somme & \casevide & \casevide \\
		Suites géométriques (termes positifs) : moyenne géométrique, terme général, somme & \casevide & \casevide \\
		Reconnaissance et détermination de la raison d’une suite arithmétique/géométrique & \casevide & \casevide \\
		Calcul de somme de termes consécutifs & \casevide & \casevide \\
		\hline
		
		% --- Fonctions exponentielles ---
		\multicolumn{3}{|L{14cm}|}{\textbf{Fonctions exponentielles}} \\
		Définition de $x \mapsto a^x$ comme prolongement des suites géométriques & \casevide & \casevide \\
		Sens de variation selon $a$, allure de la courbe & \casevide & \casevide \\
		Propriétés algébriques : $a^{x+y}$, $a^{x-y}$, $a^{nx}$ & \casevide & \casevide \\
		Taux d’évolution moyen & \casevide & \casevide \\
		\hline
		
		% --- Fonction logarithme décimal ---
		\multicolumn{3}{|L{14cm}|}{\textbf{Fonction logarithme décimal}} \\
		Définition comme solution de $10^x=b$, notation $\log$ & \casevide & \casevide \\
		Sens de variation & \casevide & \casevide \\
		Propriétés algébriques : $\log(ab)$, $\log(a^n)$, $\log(a/b)$ & \casevide & \casevide \\
		Résolution d’équations/inéquations $a^x=b$, $x^a=b$ & \casevide & \casevide \\
		\hline
		
		% --- Fonction inverse ---
		\multicolumn{3}{|L{14cm}|}{\textbf{Fonction inverse}} \\
		Comportement aux bornes, dérivée, sens de variation & \casevide & \casevide \\
		Représentation graphique, asymptotes (intuitif) & \casevide & \casevide \\
		Étude de fonctions du type $x \mapsto$ combinaison linéaire de $1/x$ et polynômes degré $\le 3$ & \casevide & \casevide \\
		\hline
		
		% --- Statistique à deux variables quantitatives ---
		\multicolumn{3}{|L{14cm}|}{\textbf{Statistique – séries à deux variables}} \\
		Nuage de points, ajustement affine (graphique et moindres carrés) & \casevide & \casevide \\
		Utilisation d’un ajustement pour interpoler/extrapoler & \casevide & \casevide \\
		Changement de variable pour linéariser & \casevide & \casevide \\
		\hline
		
		% --- Probabilités conditionnelles ---
		\multicolumn{3}{|L{14cm}|}{\textbf{Probabilités conditionnelles}} \\
		Conditionnement, indépendance de deux événements & \casevide & \casevide \\
		Formule des probabilités totales (partition en 2 ou 3 événements) & \casevide & \casevide \\
		Arbres de probabilités & \casevide & \casevide \\
		\hline
		
		% --- Variables aléatoires discrètes – Loi binomiale ---
		\multicolumn{3}{|L{14cm}|}{\textbf{Loi binomiale}} \\
		Coefficients binomiaux, triangle de Pascal, formule de Pascal & \casevide & \casevide \\
		Définition de la loi binomiale $\mathcal{B}(n,p)$, espérance (admise) & \casevide & \casevide \\
		Calcul de probabilités $P(X=k)$ pour des cas particuliers et avec coefficients & \casevide & \casevide \\
		\hline
	\end{longtable}
	
	\section{Calendrier et programmation harmonieuse}
	
	\subsection{Calendrier retenu (zone C, jours ouvrés)}
	Les vacances et jours fériés ont été pris en compte. Les semaines de cours sont les suivantes :
	
	\begin{center}
		\begin{tabular}{|C{3.5cm}|C{5cm}|C{2.5cm}|}
			\hline
			\textbf{Période} & \textbf{Dates} & \textbf{Nb de semaines} \\
			\hline
			Rentrée – Toussaint & 1er sept. – 16 oct. 2026 & 7 \\
			Toussaint – Noël & 2 nov. – 18 déc. 2026 & 7 \\
			Noël – Hiver & 4 janv. – 5 fév. 2027 & 5 \\
			Hiver – Printemps & 22 fév. – 2 avr. 2027 & 6 \\
			Printemps – Fin mai & 19 avr. – 31 mai 2027 & 6 \\
			\hline
			\textbf{Total} & & \textbf{31 semaines} \\
			\hline
		\end{tabular}
	\end{center}
	
	\textit{Jours fériés (tombant un jour de cours) exclus :} 1er janvier 2027 (vendredi), 29 mars 2027 (lundi de Pâques), 6 mai 2027 (Ascension, jeudi), 17 mai 2027 (Pentecôte, lundi). Les 1er et 8 mai 2027 tombent un samedi.
	
	\subsection{Programmation par trimestre (approche spiralée)}
	Les tableaux suivants précisent pour chaque semaine :
	\begin{itemize}
		\item les contenus nouveaux,
		\item les \textbf{points ciblés dans l’évaluation formative} (5 à 10 min en début de séance) incluant systématiquement une question Python,
		\item des remarques utiles.
	\end{itemize}
	Les évaluations formatives mêlent mathématiques et programmation, avec réactivation des notions antérieures.
	
	% --- 1er trimestre ---
	\newpage
	\subsubsection{1\textsuperscript{er} trimestre (septembre – novembre 2026) – 14 semaines}
	
	\rowcolors{2}{gray!8}{white}
	\begin{longtable}{|C{0.8cm}|L{4.0cm}|L{5.8cm}|L{2.5cm}|}
		\hline
		\textbf{Sem.} & \textbf{Contenus nouveaux} & \textbf{Évaluation formative -- points évalués (dont Python)} & \textbf{Remarques} \\
		\hline
		\endhead
		\hline
		\multicolumn{4}{r}{Suite de la page suivante} \\
		\endfoot
		\hline
		\endlastfoot
		
		1-2 & Suites : définition, suites arithmétiques (terme général, moyenne, somme). & Calculs de puissances, racines. Notion de suite (1ère). \textbf{Python} : boucle \texttt{for} pour calculer des termes et la somme. & Introduction des modèles discrets. \\
		\hline
		3-4 & Suites géométriques (termes positifs) : terme général, moyenne géométrique, somme. & \textbf{Spirale :} suites arithmétiques (S1-2). \textbf{Python} : fonction qui calcule la somme des termes d’une suite géométrique. & Lien avec les intérêts composés. \\
		\hline
		5-6 & Fonctions exponentielles $x \mapsto a^x$ : définition, sens de variation, propriétés algébriques. & \textbf{Spirale :} suites géométriques (S3-4). \textbf{Python} : tracé de courbes exponentielles avec \texttt{matplotlib}. & Modèles continus. \\
		\hline
		7-8 & Logarithme décimal : définition, sens de variation, propriétés algébriques. & \textbf{Spirale :} exponentielles (S5-6). \textbf{Python} : résolution d’équations $a^x=b$ par balayage. & Applications (échelles, pH). \\
		\hline
		9-10 & Résolution d’équations/inéquations avec logarithme (type $a^x=b$, $x^a=b$). & \textbf{Spirale :} logarithme (S7-8). \textbf{Python} : programme de recherche de seuil. & Calcul de taux moyens. \\
		\hline
		11-12 & Fonction inverse : dérivée, sens de variation, asymptotes (intuitif). Étude de fonctions du type $ax + b + c/x$. & \textbf{Spirale :} dérivation (1ère). \textbf{Python} : tracé de la courbe et de ses asymptotes. & Applications économiques (coût moyen). \\
		\hline
		13-14 & Synthèse sur les fonctions (exponentielle, logarithme, inverse). Problèmes modélisant des évolutions. & \textbf{Bilan T1 :} suites, exponentielles, log, inverse. \textbf{Python} : résolution d’équations mêlant ces fonctions. & Préparation au contrôle. \\
		\hline
		\multicolumn{4}{l}{} \\
		\hline
		\multicolumn{4}{l}{\textbf{Contrôles continus :} semaines 5, 10, 14 (avec exercice Python). \textbf{Formatives :} chaque semaine.} \\
		\hline
	\end{longtable}
	
	% --- 2e trimestre ---
	\newpage
	\subsubsection{2\textsuperscript{e} trimestre (novembre 2026 – février 2027) – 12 semaines}
	
	\rowcolors{2}{gray!8}{white}
	\begin{longtable}{|C{0.8cm}|L{4.0cm}|L{5.8cm}|L{2.5cm}|}
		\hline
		\textbf{Sem.} & \textbf{Contenus nouveaux} & \textbf{Évaluation formative -- points évalués (dont Python)} & \textbf{Remarques} \\
		\hline
		\endhead
		\hline
		\multicolumn{4}{r}{Suite de la page suivante} \\
		\endfoot
		\hline
		\endlastfoot
		
		15-16 & Statistique à deux variables : nuage de points, ajustement affine « au jugé ». & \textbf{Spirale :} représentations graphiques (automatismes). \textbf{Python} : tracé d’un nuage de points. & Contextes réels. \\
		\hline
		17-18 & Méthode des moindres carrés, droite de régression, coefficient de corrélation. & \textbf{Spirale :} ajustement (S15-16). \textbf{Python} : calcul de la somme des carrés des résidus. & Utilisation de la calculatrice. \\
		\hline
		19-20 & Changement de variable pour linéariser (ex. $y=a\,b^x$). Interpolation/extrapolation. & \textbf{Spirale :} logarithme (S7-8) pour linéariser. \textbf{Python} : régression sur des données transformées. & \textbf{PBL 1 (2h)} : Étude d’une série chronologique (ex. température). \\
		\hline
		21-22 & Probabilités conditionnelles : définitions, arbres, indépendance. & \textbf{Spirale :} probabilités de base (1ère). \textbf{Python} : simulation de tirages conditionnels. & Tests de dépistage. \\
		\hline
		23-24 & Formule des probabilités totales (partition en 2 ou 3 événements). & \textbf{Spirale :} conditionnement (S21-22). \textbf{Python} : calcul de probabilités avec arbre. & Applications en gestion. \\
		\hline
		25-26 & Loi binomiale : coefficients binomiaux, triangle de Pascal. & \textbf{Spirale :} probabilités totales (S23-24). \textbf{Python} : génération du triangle de Pascal. & Dénombrement. \\
		\hline
		\multicolumn{4}{l}{} \\
		\hline
		\multicolumn{4}{l}{\textbf{Contrôles continus :} semaines 18, 22, 26 (avec exercice Python). \textbf{Formatives :} chaque semaine.} \\
		\hline
	\end{longtable}
	
	% --- 3e trimestre ---
	\newpage
	\subsubsection{3\textsuperscript{e} trimestre (février – mai 2027) – 12 semaines}
	
	\rowcolors{2}{gray!8}{white}
	\begin{longtable}{|C{0.8cm}|L{4.0cm}|L{5.8cm}|L{2.5cm}|}
		\hline
		\textbf{Sem.} & \textbf{Contenus nouveaux} & \textbf{Évaluation formative -- points évalués (dont Python)} & \textbf{Remarques} \\
		\hline
		\endhead
		\hline
		\multicolumn{4}{r}{Suite de la page suivante} \\
		\endfoot
		\hline
		\endlastfoot
		
		27-28 & Loi binomiale : définition, calcul de probabilités, espérance (admise). & \textbf{Spirale :} coefficients binomiaux (S25-26). \textbf{Python} : simulation de la loi binomiale et histogramme. & Schéma de Bernoulli. \\
		\hline
		29-30 & Problèmes d’échantillonnage et de tests simples (sans intervalle de fluctuation). & \textbf{Spirale :} loi binomiale (S27-28). \textbf{Python} : simulation de N échantillons. & \textbf{PBL 2 (2h)} : Simulation d’un sondage. \\
		\hline
		31-32 & Synthèse probabilités : conditionnelles, totale, binomiale. & \textbf{Bilan probabilités}. \textbf{Python} : calcul automatisé de probabilités. & Préparation au contrôle. \\
		\hline
		33-34 & Réinvestissement des fonctions : optimisation avec la fonction inverse (coût moyen), problèmes de seuil avec exponentielle/logarithme. & \textbf{Spirale :} toutes les fonctions (T1). \textbf{Python} : recherche d’extremums par balayage. & \textbf{PBL 3 (2h)} : Optimisation d’un coût de production. \\
		\hline
		35-36 & Synthèse analyse et statistiques : problèmes mixtes. & \textbf{Bilan analyse + stats}. \textbf{Python} : projet complet de modélisation. & Interdisciplinarité. \\
		\hline
		37-38 & Révisions et évaluations finales. & \textbf{Bilan général :} tous les thèmes + Python. & Préparation épreuves. \\
		\hline
		\multicolumn{4}{l}{} \\
		\hline
		\multicolumn{4}{l}{\textbf{Contrôles continus :} semaines 30, 34, 38 (avec exercice Python). \textbf{Formatives :} chaque semaine.} \\
		\hline
	\end{longtable}
	
	\newpage
	
	\section{Évaluations}
	
	\subsection{Évaluations formatives (hebdomadaires)}
	Chaque semaine, une courte activité (5–10 min) mêle mathématiques et Python. Les points évalués sont détaillés dans les tableaux ci-dessus. Ces évaluations ne sont pas notées, mais elles permettent une auto-évaluation et un suivi des compétences algorithmiques.
	
	\subsection{Évaluations sommatives (contrôles continus)}
	Au moins trois contrôles par trimestre, d’une durée de 1 h chacun.  
	Chaque contrôle comporte :
	\begin{itemize}
		\item une partie sans calculatrice (calculs, raisonnement),
		\item une partie avec calculatrice / Python (exercice de programmation : écriture de fonction, complétion de code, simulation, représentation graphique, etc.).
	\end{itemize}
	
	\begin{center}
		\begin{tabular}{|C{3.5cm}|C{3.5cm}|C{3.5cm}|C{3.5cm}|}
			\hline
			\textbf{Trimestre} & \textbf{Contrôle 1} & \textbf{Contrôle 2} & \textbf{Contrôle 3} \\
			\hline
			1\textsuperscript{er} (sem. 5, 10, 14) & Suites (arithm./géom.) \newline + Python (calculs) & Exponentielle, logarithme \newline + Python (résolution) & Fonction inverse, synthèse \newline + Python (tracés) \\
			\hline
			2\textsuperscript{e} (sem. 18, 22, 26) & Statistique à 2 variables (régression) \newline + Python (nuage, résidus) & Probabilités conditionnelles \newline + Python (arbres) & Coefficients binomiaux \newline + Python (triangle) \\
			\hline
			3\textsuperscript{e} (sem. 30, 34, 38) & Loi binomiale (calculs, espérance) \newline + Python (simulation) & Optimisation, problèmes de seuil \newline + Python (balayage) & Synthèse (tous thèmes) \newline + Python (projet complet) \\
			\hline
		\end{tabular}
	\end{center}
	
	\subsection{Problèmes basés sur l’apprentissage (PBL) – 2 h chacun}
	Trois PBL sont répartis sur l’année, en séances de 2 h (travail en groupes, compte rendu écrit et code Python rendu) :
	
	\begin{enumerate}
		\item \textbf{PBL 1 (semaine 20)} : \emph{Étude d’une série chronologique}. À partir de données réelles (température, population, etc.), représenter le nuage, proposer un ajustement affine ou exponentiel (linéarisation par log), interpoler/extrapoler. \textbf{Python} : tracé, régression, transformation logarithmique.
		\item \textbf{PBL 2 (semaine 30)} : \emph{Simulation d’un sondage}. Simuler un échantillon de taille $n$ selon une loi de Bernoulli, calculer la fréquence observée, la comparer avec la probabilité théorique. \textbf{Python} : simulation, histogramme, loi binomiale.
		\item \textbf{PBL 3 (semaine 34)} : \emph{Optimisation d’un coût de production}. Étudier un coût moyen $C_m(x) = ax + b + c/x$, déterminer le minimum en utilisant la fonction inverse et la dérivée. \textbf{Python} : recherche de l’extremum par balayage et comparaison avec le calcul exact.
	\end{enumerate}
	
	\section{Devoirs maison (DM) -- progression, Python et utilisation de l’IA}
	
	Afin de favoriser une acquisition à long terme, un devoir maison est distribué \textbf{toutes les deux semaines} (18 DM sur l’année).  
	Chaque DM porte sur l’ensemble du programme traité depuis le 1er septembre, et comporte :
	\begin{itemize}
		\item des exercices de calcul et de raisonnement,
		\item un exercice de programmation Python,
		\item une mission spécifique d’utilisation de l’IA (consultation, débuggage, confrontation de méthodes ou critique de réponse).
	\end{itemize}
	
	L’élève doit joindre à sa copie un \textbf{journal de bord} retraçant ses échanges avec l’IA et les décisions prises à la suite de ces échanges.
	
	\begin{center}
		\small
		\rowcolors{2}{gray!8}{white}
		\begin{longtable}{|C{0.8cm}|C{1cm}|L{6.8cm}|L{6.8cm}|}
			\hline
			\textbf{DM} & \textbf{Sem.} & \textbf{Thèmes évalués (mathématiques + Python)} & \textbf{Mission IA spécifique} \\
			\hline
			\endhead
			\hline
			\multicolumn{4}{r}{Suite à la page suivante} \\
			\endfoot
			\hline
			\endlastfoot
			
			1 & 3 & Suites arithmétiques : terme général, somme. Python : boucle de calcul. & Demander à l’IA d’expliquer la différence entre une suite arithmétique et une suite géométrique avec un exemple concret. \\
			\hline
			2 & 5 & Suites géométriques, moyenne géométrique. Python : calcul de la somme. & Utiliser l’IA pour générer un problème de placement à intérêts composés et le résoudre. \\
			\hline
			3 & 7 & Fonctions exponentielles : propriétés, sens de variation. Python : tracé de courbes. & Demander à l’IA de justifier le sens de variation selon $a$ et de vérifier avec Python. \\
			\hline
			4 & 9 & Logarithme décimal : définition, propriétés. Python : résolution d’équations. & Utiliser l’IA pour expliquer la notion d’ordre de grandeur via le log. \\
			\hline
			5 & 11 & Résolution d’équations $a^x=b$, $x^a=b$. Python : recherche de seuil. & Demander à l’IA de proposer une méthode alternative de résolution. \\
			\hline
			6 & 13 & Fonction inverse : dérivée, asymptotes. Python : tracé. & Utiliser l’IA pour expliquer le comportement aux bornes et illustrer. \\
			\hline
			7 & 16 & Statistique : nuage de points, ajustement au jugé. Python : tracé. & Demander à l’IA de suggérer un changement de variable pour linéariser un nuage donné. \\
			\hline
			8 & 18 & Moindres carrés, droite de régression. Python : calcul des résidus. & Utiliser l’IA pour vérifier le calcul de la somme des carrés et interpréter le coefficient de corrélation. \\
			\hline
			9 & 20 & Changement de variable, interpolation/extrapolation. Python : régression sur données transformées. & Demander à l’IA de critiquer la pertinence d’une extrapolation lointaine. \\
			\hline
			10 & 22 & Probabilités conditionnelles, arbres. Python : simulation de tirage. & Utiliser l’IA pour construire un arbre de probabilités à partir d’un énoncé et vérifier les résultats. \\
			\hline
			11 & 24 & Formule des probabilités totales. Python : calcul automatisé. & Demander à l’IA d’expliquer la notion de partition avec un exemple simple. \\
			\hline
			12 & 26 & Coefficients binomiaux, triangle de Pascal. Python : génération du triangle. & Utiliser l’IA pour démontrer la formule de Pascal par un raisonnement combinatoire. \\
			\hline
			13 & 28 & Loi binomiale : définition, calcul de probabilités. Python : simulation. & Demander à l’IA de simuler un schéma de Bernoulli et de comparer avec la théorie. \\
			\hline
			14 & 30 & Espérance de la loi binomiale. Python : calcul d’espérance. & Utiliser l’IA pour interpréter l’espérance dans un contexte concret (ex. nombre de succès attendus). \\
			\hline
			15 & 32 & Échantillonnage, fluctuation. Python : simulation d’échantillons. & Demander à l’IA de discuter de la représentativité d’un échantillon et des biais possibles. \\
			\hline
			16 & 34 & Optimisation avec la fonction inverse (coût moyen). Python : recherche d’extremum. & Utiliser l’IA pour formaliser le problème et proposer une résolution algébrique. \\
			\hline
			17 & 36 & Synthèse : problèmes mêlant fonctions, statistiques, probabilités. Python : projet complet. & Demander à l’IA un retour sur la structure du code et proposer des améliorations. \\
			\hline
			18 & 38 & Bilan général. Python : projet complet. & Demander à l’IA de suggérer des extensions du problème traité. \\
			\hline
		\end{longtable}
	\end{center}
	
	\textbf{Remarques :}
	\begin{itemize}
		\item Les DM sont à rendre une semaine après la distribution.
		\item La mission IA doit être clairement identifiée dans le journal de bord (copie d’écran ou retranscription des questions/réponses, suivie de l’analyse personnelle de l’élève).
		\item La correction est effectuée en classe, avec une attention particulière aux erreurs fréquentes, aux méthodes de résolution et aux bonnes pratiques d’utilisation de l’IA.
	\end{itemize}
	
	\section{Conseils pour la mise en œuvre}
	\begin{itemize}
		\item \textbf{Rituels} : consacrer 5 à 10 minutes par séance à des exercices de calcul (mental, littéral) et à de courts extraits de code Python (lecture, complétion, débogage).
		\item \textbf{Travail en îlots} : favoriser les échanges entre élèves lors des phases de recherche et des PBL.
		\item \textbf{Numérique} : utiliser Python régulièrement (au moins une séance sur deux). Insister sur la programmation modulaire (fonctions) et la documentation du code.
		\item \textbf{IA éthique} : former les élèves à formuler des requêtes précises et à évaluer de manière critique les réponses de l’IA (vérification des calculs, cohérence, limites). Insister sur le fait que l’IA est un \emph{outil d’aide}, pas un substitut au raisonnement personnel.
		\item \textbf{Différenciation} : proposer des exercices de difficulté progressive et des approfondissements pour les élèves plus à l’aise (notamment en programmation et en interaction avec l’IA).
		\item \textbf{Trace écrite} : veiller à ce que chaque élève dispose d’un cahier de cours structuré (définitions, propriétés, exemples) et d’un cahier de programmation regroupant les fonctions et scripts clés.
	\end{itemize}
	
	\section*{Conclusion}
	Cette programmation respecte les 31 semaines de cours disponibles et couvre l’intégralité du programme de terminale technologique (tronc commun), avec une intégration systématique de Python et de l’utilisation critique de l’IA dans les apprentissages et les évaluations. La progressivité spiralaire, appliquée aussi aux compétences algorithmiques et numériques, garantit que les notions mathématiques, la programmation et le jugement critique sont solidement maîtrisés en fin d’année. Les évaluations variées (formatives, sommatives, PBL, DM) permettent de suivre la progression de chaque élève et de développer ses compétences de modélisation, de raisonnement, de communication, de codage et de collaboration avec les outils d’intelligence artificielle.
	
\end{document}